\documentclass[11pt]{article}
\usepackage[margin=1in]{geometry}
\emergencystretch=2em
\usepackage{amsmath,amssymb,amsthm}
\usepackage{hyperref}

\newtheorem{theorem}{Theorem}
\newtheorem{lemma}{Lemma}
\newtheorem{remark}{Remark}

\newcommand{\R}{\mathbb{R}}
\newcommand{\Acal}{\mathcal{A}}
\newcommand{\G}{\mathsf{G}_{\mathcal A}}

\title{Finite-Dimensional Positive Case for the Aronszajn-Null Game}
\author{Run \texttt{fa\_banach\_001}, agent \texttt{agent\_lane\_00}}
\date{2026-06-15}

\begin{document}
\maketitle

\section*{Source and Open Question}

The source paper is Jakub Duda and Boaz Tsaban, \emph{Games in Banach spaces},
arXiv:math/0601556.  For a separable Banach space \(X\), and a non-zero
direction \(x\in X\), let \(\mathcal A(x)\) be the Borel sets whose
intersection with every line parallel to \(x\) has one-dimensional Lebesgue
measure zero.  A Borel set \(A\subset X\) is Aronszajn-null if for every dense
sequence \((x_n)\) in \(X\) there are \(A_n\in\mathcal A(x_n)\) such that
\[
A\subset \bigcup_n A_n.
\]

Duda--Tsaban define the game \(\G\): in inning \(n\), Player I plays a direction
\(x_n\) so that the resulting sequence is dense, Player II answers with
\(A_n\in\mathcal A(x_n)\), and Player II wins if \(A\subset\bigcup_n A_n\).
They prove that if Player I has no winning strategy, then \(A\) is
Aronszajn-null, and state the converse as open:
\[
A\in\mathcal A \quad\Longrightarrow\quad \text{Player I has no winning
strategy in }\G.
\]

\section*{Partial Result}

\begin{theorem}
Let \(X\) be a finite-dimensional real Banach space and let \(A\subset X\) be a
Borel Aronszajn-null set. Then Player II has a winning strategy in the
Aronszajn-null game \(\G\) played on \(A\).
\end{theorem}

\section*{Proof Intuition}

In finite dimensions, Aronszajn-null Borel sets are Lebesgue null. Suppose
Player I has already played independent directions spanning a subspace \(S\).
The part of \(A\) not yet covered can be over-covered by the inverse image of a
null Borel set in the quotient \(X/S\). When Player I next plays a direction
whose image in \(X/S\) is non-zero, Fubini splits that quotient-null set into
the part with null sections parallel to the new direction and a smaller
quotient-null residual. Pulling the first part back to \(X\) is a legal move
for Player II. Since a dense sequence in a finite-dimensional space spans the
whole space after finitely many independent moves, this process exhausts the
residual.

\section*{Proof}

We use Lebesgue measure on finite-dimensional quotients of \(X\), normalized
arbitrarily; nullness is independent of the normalization.

First note that \(A\) is Lebesgue null. Indeed, since \(A\) is Aronszajn-null,
take any dense sequence containing a basis of \(X\). Then \(A\) is covered by
countably many sets \(B_n\in\mathcal A(x_n)\). For any non-zero \(x\), each
Borel \(B\in\mathcal A(x)\) is Lebesgue null by Fubini over the decomposition
of \(X\) into lines parallel to \(x\). Hence \(A\) is null.

We describe a strategy for Player II. More generally, maintain the following
state after some innings. Let \(S\subset X\) be the span of the independent
directions already accepted, let \(\pi:X\to Q=X/S\) be the quotient map, and
let \(R\subset Q\) be a Borel null set such that the part of \(A\) not yet
covered by Player II's previous moves is contained in \(\pi^{-1}(R)\). Initially
\(S=\{0\}\), \(Q=X\), \(R=A\), and \(\pi\) is the identity.

Suppose Player I plays \(v\). If the image \(w=\pi(v)\) is zero, Player II plays
\(\emptyset\) and leaves the state unchanged. Otherwise let
\(q:Q\to Q/\R w\) be the quotient map. By the standard Fubini theorem for Borel
sets, the section-measure function
\[
z\mapsto m_1\bigl(R\cap q^{-1}(z)\bigr)
\]
is measurable and
\[
N=\{z\in Q/\R w: m_1(R\cap q^{-1}(z))>0\}
\]
is a null Borel set in \(Q/\R w\), because \(R\) is null in \(Q\). Set
\[
G=R\setminus q^{-1}(N).
\]
Then every line in \(Q\) parallel to \(w\) meets \(G\) in a one-dimensional null
set. Player II plays
\[
B=\pi^{-1}(G).
\]
This is a legal move: if \(L\) is any line in \(X\) parallel to \(v\), then
\(\pi(L)\) is a line in \(Q\) parallel to \(w\), and \(L\cap B\) is identified
with a subset of \(\pi(L)\cap G\), hence has one-dimensional measure zero.
Thus \(B\in\mathcal A(v)\).

The uncovered part of \(A\) is now contained in
\[
\pi^{-1}(R\setminus G)\subset \pi^{-1}(q^{-1}(N)).
\]
Replace \(S\) by \(S+\R v\), replace \(Q\) by \(Q/\R w\), replace \(\pi\) by
\(q\circ\pi\), and replace \(R\) by \(N\). The invariant is preserved and the
dimension of the quotient has dropped by one.

A legal play by Player I must be dense in \(X\), so its played directions span
\(X\). Therefore Player I eventually plays \(\dim X\) linearly independent
directions. After the last such move the quotient has dimension zero. In the
one-dimensional final quotient step, the null set \(R\) has zero
one-dimensional measure, so the positive-section set \(N\) is empty. Hence the
remaining residual is empty. Thus \(A\) is covered by Player II's moves, and
the described strategy wins.
\qed

\section*{Verification Notes}

The proof uses only standard finite-dimensional Fubini disintegration for
Borel sets under a linear quotient map. The moves played by Player II need not
be subsets of \(A\); they only need to be Borel members of the corresponding
\(\mathcal A(v)\), as in the source game definition.

If one permits Player I to play the zero vector despite \(\mathcal A(0)\) not
being defined in the source paper, the strategy treats that inning as
irrelevant and plays \(\emptyset\). Equivalently, one may read the game as
requiring non-zero directions, which is the only interpretation under which
Player II's legal response is defined.

\section*{Limitations and Novelty Check}

This packet proves only the finite-dimensional case. It does not address the
source conjecture for infinite-dimensional separable Banach spaces, nor the
other diagram-reversal questions in the paper.

Cheap run indexes had no previous packet for arXiv:0601556. Local corpus
searches for ``Aronszajn-null game'' and exact phrases from the source
conjecture found only the source paper. A bounded web search for the same
phrases likewise found the source arXiv page and no later answer. Novelty
confidence is moderate: the finite-dimensional argument is elementary and may
be folklore, but it is not recorded in the run memory.

\section*{References}

\begin{itemize}
\item J. Duda and B. Tsaban, \emph{Games in Banach spaces},
arXiv:math/0601556.
\item Y. Benyamini and J. Lindenstrauss, \emph{Geometric Nonlinear Functional
Analysis, Volume 1}, AMS Colloquium Publications 48, 2000.
\end{itemize}

\end{document}
